\section{Decision Inference}
\sectionkicker{INTERFACES AND WEIGHTS}
\begin{wideflow}
{\Large\bfseries 判断支援の3界面――OllamaのAPI、Clefの重みとAPI、Strands Deciderのpointer head}\par
\smallskip
{\color{SurveyMuted}9月29日のインターフェースと10月1日の重み・モデルを混ぜない。入力と出力、測定条件、ライセンスを面ごとに読む。}\par
\end{wideflow}

Ollamaは9月29日、v0.35で判断支援モデル向けのインターフェースを開けた\cite{ollama}。/v1/systemoneというAPIに、Nimble 9BとTev1の4B・0.8Bモデルが並ぶ。ブログ本文は通読済みの範囲の話である。91ミリ秒という数値は単一事例の例示であって、普遍的な性能ではない。ここは新しいモデルの登場ではなくインターフェースの話として読む。レジストリやAPI文書の差分は積み残しである。

Cloudflareは10月1日15時34分02秒（UTC）にClefをリリースした\cite{clef}。重みはApache-2.0ライセンスで公開され、Jev API、6万4千トークンのコンテキストと画像入力、入力0.24ドル・出力0.09ドルの料金が並ぶ。入力は意思決定タスクの記述と画像、出力は判断結果と根拠という構成である。Jev Indexの先行結果やレイテンシはベンダーが主催するデモ環境での計測であり、独立した再現ではない。重み配布ページの存在は確認したが、重みファイルの取得まではしていない。OllamaのインターフェースやStrandsのモデルと数値・レイテンシ・ライセンスを混ぜない。

Strandsは10月1日、Decider 2Bをリリースした\cite{strands}。Qwen3.5-2Bベースにpointer headとr16 LoRAを添えた構成で、ベンダー帰属つきのJevBench結果とレイテンシ、Apache-2.0ライセンスの三重確認が示される。重み配布ページの版はv19で固定する。先頭表示が古い版の様相を見せる配布ページの癖があるため、版の管理は欠かさない。ハイブリッドエージェント向けという位置づけはベンダーの助言として読む。

時系列の整理として、9月29日はOllamaのインターフェース公開、10月1日はClefの重み・API公開とStrands Decider 2Bのモデル公開である。入力・出力・測定条件・ライセンスの4点を面ごとに切り分けると、3者は判断支援という同じ層にいながら技術面が重ならないことがわかる。

\begin{claimboundary}[資料と評価の境界]
インターフェースの話と重みの話とモデルの話は別の面であり、数値や許諾や待ち時間を混ぜない。91ミリ秒は一点の例示であり、普遍の速さではない。Jev指数や待ち時間はベンダーが主催するデモの計測であり、独立の再現ではない。重みの中身の取得や文書の差分は積み残しである。
\end{claimboundary}
